\documentclass[a4,11pt]{aleph-notas}
% Se puede ver la documentación aquí:
% https://github.com/alephsub0/LaTeX_aleph-notas

% -- Paquetes adicionales
\usepackage{enumitem}
\usepackage{url}
\usepackage{array}
\usepackage{booktabs}
\usepackage{longtable}
\usepackage{aleph-comandos}
\hypersetup{
    urlcolor=blue,
    linkcolor=blue,
}


% -- Datos
\institucion{Facultad de Ciencias Exactas, Naturales y Ambientales}
\carrera{Catálogo STEM}
\asignatura{Matemática III}
\tema{Plan tema 04: Límites y diferenciabilidad}
\autor{Andrés Merino}
\fecha{Periodo 2026-1}

\logouno[0.16\textwidth]{Logos/logoPUCE_04_ac}
\definecolor{colortext}{HTML}{1957d1}
\definecolor{colordef}{HTML}{1957d1}
\fuente{montserrat}


\begin{document}

\encabezado

%%%%%%%%%%%%%%%%%%%%%%%%%%%%%%%%%%%%%%%%
\section*{Resultado de Aprendizaje}
%%%%%%%%%%%%%%%%%%%%%%%%%%%%%%%%%%%%%%%%

%%%%%%%%%%%%%%%%%%%%%%%%%%%%%%%%%%%%%%%%
\subsection*{RdA de la asignatura:}
\begin{itemize}[leftmargin=*]
    \item \textbf{RdA 1:} Identificar los conceptos, teoremas y propiedades de funciones vectoriales, derivadas parciales e integral vectorial.
    \item \textbf{RdA 2:} Calcular derivadas parciales e integrales de funciones vectoriales con el apoyo de asistentes computacionales.
    % \item \textbf{RdA 3:} Aplicar distintos tópicos del Cálculo Vectorial para la solución de problemas reales en diferentes contextos.
\end{itemize}

%%%%%%%%%%%%%%%%%%%%%%%%%%%%%%%%%%%%%%%%
\subsection*{RdA del tema:}
\begin{itemize}[leftmargin=*]
    \item Calcular límites de campos escalares de varias variables.
    \item Determinar la no existencia de un límite mediante trayectorias.
    \item Verificar la continuidad de un campo escalar en un punto.
    \item Calcular derivadas parciales y direccionales de un campo escalar.
    \item Determinar si un campo escalar es diferenciable en un punto y calcular su diferencial.
    \item Aplicar la regla de la cadena para campos escalares compuestos con trayectorias.
\end{itemize}

%%%%%%%%%%%%%%%%%%%%%%%%%%%%%%%%%%%%%%%%
\section*{Introducción}
%%%%%%%%%%%%%%%%%%%%%%%%%%%%%%%%%%%%%%%%

\paragraph{Pregunta inicial:}
Si la temperatura de un cuarto depende de la posición $(x,y)$ donde estés parado, ¿cómo describirías que la temperatura «no cambia abruptamente» al moverte? ¿Importa desde qué dirección te acerques a un punto?

%%%%%%%%%%%%%%%%%%%%%%%%%%%%%%%%%%%%%%%%
\section*{Desarrollo}
%%%%%%%%%%%%%%%%%%%%%%%%%%%%%%%%%%%%%%%%

%%%%%%%%%%%%%%%%%%%%%%%%%%%%%%%%%%%%%%%%
\subsection*{Actividad 1: Límites y continuidad}
Se define el límite de un campo escalar de varias variables y se discuten las estrategias para calcularlo o demostrar su no existencia: propiedades algebraicas, cambio de variable, trayectorias y el teorema del emparedado.

\paragraph{¿Cómo lo haremos?}
\begin{itemize}[leftmargin=*]
    \item \textbf{Motivación:} retomar la pregunta inicial y comparar con límites en una variable; enfatizar que en varias variables el límite debe coincidir por \emph{todas} las trayectorias.
    \item \textbf{Clase magistral:} se definen límite, límite por trayectorias, continuidad y los teoremas de propiedades, cambio de variable y emparedado. Se usa el resumen \href{https://andres-merino.github.io/MatematicaIII-05-N0051/01Resumen/Resumen04.pdf}{Resumen04.pdf}.
    \item \textbf{Resolución de ejercicios:} los estudiantes calcularán límites usando distintas estrategias y verificarán la no existencia por trayectorias, usando \href{https://andres-merino.github.io/MatematicaIII-05-N0051/04Ejercicios/Ejercicios04.pdf}{Ejercicios04.pdf}.
\end{itemize}

\paragraph{Verificación de aprendizaje:}
\begin{itemize}[leftmargin=*]
    \item ¿Por qué encontrar dos trayectorias con límites distintos prueba que el límite no existe?
    \item ¿Qué condición adicional distingue un límite que existe de la continuidad en un punto?
    \item ¿Cómo se aplica el cambio de variable para calcular $\dlim_{(x,y)\to(0,0)}\frac{\sin(x^2+y^2)}{x^2+y^2}$?
\end{itemize}

%%%%%%%%%%%%%%%%%%%%%%%%%%%%%%%%%%%%%%%%
\subsection*{Actividad 2: Diferenciabilidad}
Se introduce la derivada direccional y parcial, el gradiente y el concepto de diferenciabilidad como aproximación lineal. Se establece la jerarquía: diferenciable $\Rightarrow$ continuo, diferenciable con continuidad $\Rightarrow$ diferenciable, y se presentan contraejemplos para las implicaciones inversas.

\paragraph{¿Cómo lo haremos?}
\begin{itemize}[leftmargin=*]
    \item \textbf{Clase magistral:} se definen derivada respecto a un vector, derivadas parciales, gradiente y diferenciabilidad; se enuncia la condición suficiente (derivadas parciales continuas) y la regla de la cadena. Se usa el resumen \href{https://andres-merino.github.io/MatematicaIII-05-N0051/01Resumen/Resumen04.pdf}{Resumen04.pdf}.
    \item \textbf{Implementación computacional:} calcular gradientes y verificar diferenciabilidad con un asistente computacional para campos escalares concretos.
    \item \textbf{Resolución de ejercicios:} los estudiantes calcularán derivadas parciales, gradientes y aplicarán la regla de la cadena, usando \href{https://andres-merino.github.io/MatematicaIII-05-N0051/04Ejercicios/Ejercicios04.pdf}{Ejercicios04.pdf}.
\end{itemize}

\paragraph{Verificación de aprendizaje:}
\begin{itemize}[leftmargin=*]
    \item ¿En qué dirección unitaria crece más rápido un campo escalar en un punto dado?
    \item ¿Por qué la existencia del gradiente no garantiza la diferenciabilidad?
    \item ¿Cuál es la relación entre el diferencial $f'(a)$ y el gradiente $\nabla f(a)$?
\end{itemize}

%%%%%%%%%%%%%%%%%%%%%%%%%%%%%%%%%%%%%%%%
\section*{Cierre}
%%%%%%%%%%%%%%%%%%%%%%%%%%%%%%%%%%%%%%%%

\paragraph{Tarea:}
Resolver del libro \href{https://catalogobiblioteca.puce.edu.ec/cgi-bin/koha/opac-detail.pl?biblionumber=83405&query_desc=su%3A%22An%C3%A1lisis%20vectorial%22}{Cálculo Vectorial de Colley}: sección 2.2, ejercicios 7, 9, 11, 13, y sección 2.3, ejercicios 1--25 (impares).

\paragraph{Pregunta de investigación:}
\begin{enumerate}[leftmargin=*]
    \item ¿Cómo se usa el gradiente en algoritmos de aprendizaje automático para minimizar funciones de error?
    \item ¿Qué significa geométricamente que el gradiente sea perpendicular a las curvas de nivel de un campo escalar?
\end{enumerate}

\paragraph{Para el próximo tema:}
Ver los videos:
\begin{itemize}[leftmargin=*]
    \item \href{https://youtu.be/Ag9GBB3pPkU?si=TgUNvD-IEr8Opy6x}{¿Por qué el Gradiente se ve de esta manera?}.
    \item \href{https://youtu.be/A6FiCDoz8_4?si=J93rcxDWUy4mN7G0}{¿Qué es el Descenso del Gradiente?}.
\end{itemize}


\end{document}
